\chapter{Conclusiones y Recomendaciones}
\pagestyle{fancy}

\section{Conclusiones}
La evolución del sistema de gestión académica hacia una arquitectura orientada a microservicios permitió resolver satisfactoriamente las carencias de escalabilidad, acoplamiento y despliegue que afectaban al monolito original. La segregación de responsabilidades en dominios independientes —tales como identidad, lógica académica y notificaciones— simplificó la comprensión del sistema, redujo el riesgo de fallas en cascada y permitió adoptar tecnologías modernas adaptadas a las necesidades específicas de cada módulo.

En paralelo, la incorporación de un modelo de auto-mantenimiento demostró ser una solución pragmática e innovadora frente a la escasez de recursos humanos dedicados a la administración operativa. El sistema desarrollado no solo logró consolidar las métricas de rendimiento y los registros de error dispersos, sino que fue capaz de inferir proactivamente soluciones mediante la integración de capacidades de Inteligencia Artificial (LLMs). Las pruebas de caos demostraron que el sistema es capaz de detectar la caída de servicios, alertar inmediatamente a las partes interesadas mediante notificaciones institucionalizadas y proveer al operador un entorno de diagnóstico asistido e inmediato.

En definitiva, se comprobó empíricamente que la observabilidad distribuida, combinada con algoritmos generativos orientados a infraestructura (AIOps), acelera drásticamente el diagnóstico de fallas (Root Cause Analysis), reduciendo los tiempos de recuperación (MTTR) e incrementando la disponibilidad del sistema.

\section{Trabajo Futuro y Recomendaciones}
Para potenciar el alcance y la autonomía del modelo implementado, se sugieren las siguientes vías de desarrollo futuro. La visión a largo plazo es dotar a PAWARE de un contexto mucho más profundo y herramientas de acción avanzadas, evolucionándolo de un asistente de diagnóstico a un agente de ingeniería de software completamente autónomo:

\begin{itemize}
    \item \textbf{Integración Profunda con el Código Fuente (Control de Versiones):} Ampliar las capacidades de la IA para que pueda conectarse a los repositorios del proyecto (ej. GitHub/GitLab). Esto le permitiría, ante una alerta, clonar el código, auditar la lógica, encontrar el segmento defectuoso y sugerir —o incluso generar automáticamente— \textit{Pull Requests} con la refactorización o corrección (auto-mejora del código).
    \item \textbf{Interacción Activa con Servicios y Datos:} Evolucionar a PAWARE para que no solo consuma métricas pasivas de observabilidad, sino que pueda ejecutar consultas a través de las APIs internas o inspeccionar el estado de las bases de datos de forma segura. Si un operador consulta: \textit{``¿Cómo podemos mejorar el rendimiento del módulo de asignaturas?''}, la IA podría contrastar la telemetría con los volúmenes de datos reales y emitir recomendaciones de índices de base de datos o rediseño de queries.
    \item \textbf{Tareas Programadas y Auditorías Predictivas:} Habilitar a PAWARE para ejecutar ciclos de revisión preventiva programados (cron jobs de auditoría). En lugar de reaccionar a incidentes, la IA simularía flujos de carga internamente, evaluaría la deuda técnica o los cuellos de botella latentes, y emitiría informes periódicos de mejora continua antes de que los usuarios experimenten lentitud.
    \item \textbf{Ejecución Autónoma en Orquestadores:} Avanzar del estado de "recomendación" a la acción directa sobre la infraestructura (Kubernetes o Docker Swarm), otorgando permisos controlados a la IA para reiniciar, apagar o escalar instancias horizontalmente sin confirmación humana explícita, completando la visión de un sistema auto-reparable (Self-healing).
    \item \textbf{Modelos Específicos y Soberanía Tecnológica (SLMs):} Recopilar progresivamente las resoluciones aprobadas en la base de datos de conocimiento para entrenar modelos de lenguaje de menor escala (Small Language Models). Esto lograría que el asistente capture la terminología exclusiva de la Universidad Central de Venezuela y opere localmente, reduciendo costos y dependencia de APIs de terceros.
\end{itemize}

